\documentclass[tikz,border=10pt]{standalone}
\usepackage{tikz}
\usepackage{amsmath}
\usepackage{amssymb}
\usepackage{pgfplots}
\pgfplotsset{compat=1.18}
\usepackage{ctex}
\usetikzlibrary{calc, positioning, arrows.meta}

\begin{document}
\begin{tikzpicture}

\begin{axis}[
  width=9cm, height=6.5cm,
  xmin=0, xmax=3.4,
  ymin=-0.3, ymax=1.5,
  xtick={0, 1.5708, 3.1416},
  xticklabels={$0$, $\frac{\pi}{2}$, $\pi$},
  ytick={0, 0.5, 1.0},
  yticklabels={$0$, $0.5\|\mathbf{u}\|\|\mathbf{v}\|$, $\|\mathbf{u}\|\|\mathbf{v}\|$},
  xlabel={$\theta$},
  ylabel={$|\langle\mathbf{u},\mathbf{v}\rangle|$},
  xlabel style={font=\small},
  ylabel style={font=\small},
  tick label style={font=\small},
  axis lines=left,
  axis line style={gray!60},
  grid=major,
  grid style={gray!20, thin},
  clip=false,
]

  % 余弦绝对值曲线 |cos(x)|，纵轴单位归一化（乘以 ||u||·||v||=1）
  \addplot[
    domain=0:3.1416,
    samples=200,
    thick,
    blue!70!black,
  ] {abs(cos(deg(x)))};

  % 上界水平直线 ||u||·||v||
  \addplot[
    domain=0:3.1416,
    thick,
    red!70!black,
    dashed,
  ] {1.0} node[right, font=\small, red!80!black] {$\|\mathbf{u}\|\|\mathbf{v}\|$（上界）};

  % 等号成立标记：theta=0
  \addplot[mark=*, mark options={fill=red,draw=red,scale=1.2}]
    coordinates {(0, 1.0)};
  \node[font=\scriptsize, above right, align=center, text=red!80!black]
    at (axis cs: 0, 1.0) {等号成立\\$\theta=0$\\（共线同向）};

  % 等号成立标记：theta=pi
  \addplot[mark=*, mark options={fill=red,draw=red,scale=1.2}]
    coordinates {(3.1416, 1.0)};
  \node[font=\scriptsize, above left, align=center, text=red!80!black]
    at (axis cs: 3.1416, 1.0) {等号成立\\$\theta=\pi$\\（共线反向）};

  % theta=pi/2 处：最小值 0
  \addplot[mark=o, mark options={fill=blue!30,draw=blue!70,scale=1.1}]
    coordinates {(1.5708, 0)};
  \node[font=\scriptsize, below right, text=blue!70!black]
    at (axis cs: 1.5708, 0) {$\theta=\tfrac{\pi}{2}$\\（正交，内积为 0）};

\end{axis}

% 公式框（放在 axis 外下方）
\node[font=\small, align=center, fill=white, draw=gray!50, thin,
      rounded corners=3pt, inner sep=7pt]
  at (4.5, -1.1) {
    Cauchy-Schwarz 不等式：
    $|\langle\mathbf{u},\mathbf{v}\rangle| \le \|\mathbf{u}\|\,\|\mathbf{v}\|$\\
    等号成立 $\Leftrightarrow$ $\mathbf{u},\mathbf{v}$ 共线
  };

% 标题
\node[font=\large\bfseries] at (4.5, 5.3) {Cauchy-Schwarz 不等式的几何意义};

\end{tikzpicture}
\end{document}
